
\section{Context Description}

\subsection{General Overview}
This project involves the design of an extensive and centralized software application, intended for use in the administration of the operations of a network of private healthcare services spread over the Italian territory. This application, operating exclusively on an outpatient and ambulatory model, seeks to digitalize key administrative operations, thereby optimizing the entire experience for all relevant stakeholders involved in the healthcare sector.

At the heart of this application is the ``Examination,'' defined as a specific, scheduled transaction between a patient and a doctor. In order to support this model, the application will offer medical professionals tools for scheduling, administrators an extensive toolset for overseeing hospital operations, and patients an optimized interface for scheduling examinations and processing transactions.

The specific features and step-by-step operations for each of these stakeholders will be discussed in the objectives and requirements sections of the application.

\subsection{System Boundaries}
A major technical limit of this system is its ability to integrate with third-party insurance providers. The application will be unable to verify policies in real time, as it will not have access to any external, live databases from third-party insurance providers. The system will, therefore, rely on a simulated system for managing patient insurance information, with codes created internally. The process of reconciling insurance claims and outstanding debts owed to the hospital will, therefore, not be automated but will be internally simulated.

In addition, the fact that the platform has been specifically designed for the logistics of the scheduled outpatient examinations means that the current scope of the system has been specifically defined in terms of doctors, facility managers, and patients. There has been no specific design of the system in terms of dedicated interfaces and access levels for other professional figures in the hospital scenario, such as nurses, medical technicians, and administrative secretaries.
      
Furthermore, the administration of medications and pharmaceuticals is completely out of the scope of this project. The functionality of the application still remains completely focused on the administration and processing of medical consultations.
\newpage

\newpage